\section{Factorización de los lectores terminales mediante naturalidad, refinamiento y paso al límite}

El punto de partida material es el único constructor simultáneo que produce un
solo continuo enriquecido con cinco operaciones intrínsecas correlativas. Las
notaciones
\(\mathfrak G_{\rm sp}\), \(\mathfrak G_{\rm sol}\),
\(\mathfrak G_{\rm gau}\), \(\mathfrak G_{\rm coh}\) y
\(\mathfrak G_{\rm det}\) designan cinco vistas completas de ese mismo registro,
reproducido con sus fibras,
relaciones, números, orientación, acarreo, memoria, frontera, ruta, multisección,
supervivencia, terminales y lectores. No se toman como cinco ramas elegidas ni como
cinco datos independientes: se demuestra su compatibilidad sobre el mismo ledger,
su naturalidad bajo refinamiento, su no-vacuidad conjunta y su paso al límite, y
después se aplican, respectivamente, a los cinco problemas. Ninguna prueba de este
volumen remite a otra fuente externa para completar una definición, una identidad o un paso
lógico.

Cada cierre se expone de forma autocontenida: las definiciones, demostraciones y
compatibilidades aparecen antes de su conclusión clásica. En cada resolución
se exhiben cuatro capas distintas:
\begin{enumerate}
  \item identidad exacta en cada nivel finito;
  \item compatibilidad de refinamientos y preservación de la historia enriquecida;
  \item uniformidad y paso al límite, con su terminal conservado;
  \item reconocimiento de la conclusión clásica en el codominio publicado.
\end{enumerate}

La composición rectora de las cinco clausuras conserva primero la acción única
correlativa, forma el continuo sólo al pasar al límite y después abre una vista para
cada aplicación:
\[
\begin{aligned}
 \mathcal E_k&\xrightarrow{\ \mathcal O_k\ }\mathcal Z_k
 \xrightarrow{\ \varprojlim_k\ }
 \bigl(\mathcal C_{\rm cont}^{\rm disc},\mathcal O_\infty\bigr),\\
 \bigl(\mathcal C_{\rm cont}^{\rm disc},\mathcal O_\infty\bigr)
 &\xrightarrow{\ \operatorname{View}_T\ }\mathfrak G_T^{\rm int}
 \xrightarrow{\ \operatorname{App}_{T,d_T}\ }\mathfrak G_T(d_T),\\
 \mathfrak G_T(d_T)&
 \xrightarrow{\widetilde{\mathcal A}_{T,d_T}}\mathfrak M_T
 \xrightarrow{\mathcal P_T}\mathcal C_T,
 \qquad T\in\{\mathrm{sp,sol,gau,coh,det}\}.
\end{aligned}
\]
Aquí $\mathcal O_k$ actúa una sola vez sobre el estado completo y conserva juntas
las cinco características. $\operatorname{View}_T$ no restringe el dominio ni
selecciona historias: sólo expone las coordenadas de una aplicación dentro del valor
completo de $\mathcal O_\infty$. El dato clásico $d_T$ entra sólo en
$\operatorname{App}_{T,d_T}$, después del continuo, y no participa en la generación
de la característica intrínseca. La recta real es asimismo una salida posterior de
$\operatorname{ev}_{\mathbb R}$, no una premisa de la cadena.

Cada capítulo distingue el dominio de partida, la construcción, el certificado,
la conclusión demostrada y sus falsadores exactos. Los propietarios precedentes
construyen las cinco composiciones en sus dominios respectivos; los capítulos que
siguen despliegan esas construcciones y efectúan el reconocimiento final. Ninguna
flecha se reemplaza por una remisión, una conclusión nominal o una hipótesis clásica.
